\scanpage{115}
\mspage{2.4}
to apply proposition 2.4 b) from which we conclude that $r((US)_{\underline p}/R)=r(S/R)$.
Therefore we have that $r(S/R)=r((US)_{\underline p}/R)=r((US)_{\underline p}/U)$.

The following lemma is well-known ([12]).

Lemma 2.6: Let $R$ be a local domain (not necessarily normal) and $S$ an
$R$-algebra which is finitely generated as $R$-module. Then for each ideal of
definition $\underline q$ in $R$ we have $e_R(q)[S:R]=\sum_i e_{S_{Mi}}(qS_{Mi})[\overline{S_{Mi}}:\overline R]$ where $Mi$
runs through the maximal ideals of $S$. In particular if $S$ is a finite local
$R$-algebra then $e_{S/R}(\underline q)=[S:R]/[\overline S:\overline R]$ for each ideal of definition $\underline q$
in $R$ where $[S:R]=[S\otimes_R k:k]$ and $k$ is the quotient field of $R$.

We now show that for integrally closed local algebras the ramification
index as defined in the introduction has the essential properties one
usually expects.

Theorem 2.7: Let $R\subset S\subset T$ be integrally closed local rings which are
local algebras. Then,

a) Let $U$ be a transporter for $S$ (see end of section 1 for definition).
Then $r(S/R)=[US:U]$ and is thus an integer.

b) $r(S/R)=1$ if and only if $S$ is an unramified $R$-algebra.

c) $r(T/R)=r(T/S)r(S/R)$.

d) $r(S/R)=e(S/R)[\overline S:\overline R]_i$ where $[\overline S:\overline R]_i$ is the degree of inseparability
of the field extension $\overline S$ of $\overline R$.

\uline{Proof}: a) By proposition 1.3, we know that there always exists a
transporter $U$ for the local $R$-algebra $S$. Since $U$ is a transporter we know
that $U$ is an unramified local $R$-algebra and $US$ is a finite local $U$-algebra.
Therefore by proposition 2.5 we know that $r(S/R)=r(US/U)$. The other
property the transporter has is that $\overline{US}$ is a purely inseparable extension
of $\overline U$. Therefore $US\otimes_U US$ is a local ring and consequently $(US)_U^f=US\otimes_U US$.
\mspage{114}
